\documentclass[conference]{IEEEtran}
\usepackage{cite}
\usepackage{amsmath,amssymb,amsfonts}
\usepackage{algorithmic}
\usepackage{graphicx}
\usepackage{textcomp}
\usepackage{xcolor}
\usepackage{hyperref}
\usepackage{booktabs}
\usepackage{algorithm}
\usepackage{multirow}

\begin{document}

\title{AI-Assisted Automated Timetable Scheduling Using Constraint Satisfaction and Multi-Model LLM Fallback}

\author{
\centering
\begin{tabular}{c@{\hspace{0.6cm}}c@{\hspace{0.6cm}}c@{\hspace{0.6cm}}c}
\textbf{Akash A} &
\textbf{Sreekaram Harshith} &
\textbf{Mane Rakesh} &
\textbf{Tanveer Ahmed} \\

20221CIT0110 &
20221CIT0104 &
20221CIT0105 &
Project Guide \\[0.5cm]

\multicolumn{4}{c}{Department of Computer Science and Engineering} \\
\multicolumn{4}{c}{Presidency University, Bengaluru}
\end{tabular}
}

\maketitle

\begin{abstract}
Scheduling timetables for universities is a hard problem. Rooms clash, faculty get double-booked, student batches overlap, and the whole thing usually ends up being done by hand on a spreadsheet. This paper describes a web-based system that automates timetable generation using a Constraint Satisfaction Problem (CSP) solver with a Minimum Remaining Values (MRV) heuristic, and layers a Retrieval-Augmented Generation (RAG) chatbot on top so users can ask questions about their schedules in plain language. The CSP engine handles hard constraints (no double-booking of rooms, faculty, or batches) and optimizes for soft ones (balanced workloads, fewer gaps). The chatbot queries the live database, builds context from keyword extraction and intent classification, and routes the prompt to one of four free LLM models through a sequential fallback mechanism. We tested the system with synthetic data representing a mid-sized institution (6 departments, 40 faculty, 120 subjects, 30 rooms) and found that the generator resolves conflict-free timetables in under 30 seconds for most configurations. The RAG system answers schedule queries with database-grounded responses in 2--5 seconds. The full system runs on a free-tier deployment stack (Vercel + MongoDB Atlas + OpenRouter), making it accessible to institutions that cannot afford commercial scheduling software or paid API subscriptions.
\end{abstract}

\begin{IEEEkeywords}
Timetable Scheduling, Constraint Satisfaction Problem, Retrieval-Augmented Generation, Multi-Model LLM, Next.js, MongoDB
\end{IEEEkeywords}

\section{Introduction}

Every semester, someone at every university gets stuck with the timetable. They open a spreadsheet, pull up the list of subjects and faculty, and start slotting things in manually. Two hours later they realize Prof.\ Kumar is teaching in two rooms at 10 AM, Lab 3 is booked for three different batches on Tuesday, and the second-year students have a four-hour gap between classes. They start over.

This is not a new problem. University course timetabling has been studied since the 1960s \cite{even1975complexity}, and it is known to be NP-hard in its general form \cite{cooper1996complexity}. The search space grows combinatorially: for $n$ subjects, $m$ time slots, $r$ rooms, and $f$ faculty members, the number of possible assignments is $O((m \times r \times f)^n)$. Even a small department can produce billions of candidate schedules, most of which violate at least one constraint.

Automated scheduling tools exist, but the commercial ones are expensive and the open-source ones tend to be research prototypes that never quite work in production. Most assume you have a dedicated system administrator who understands the configuration format. None of them let a confused student type "when is my Database Systems class?" and get an answer.

We built a system that does two things: it generates timetables automatically using a CSP solver, and it lets users interact with those timetables through a conversational AI assistant. The system supports five user roles (Admin, Head of Department, Coordinator, Faculty, Student), each with a different view of the same underlying data. The AI chatbot is available on every dashboard and answers questions using actual database records, not hallucinated guesses.

We made a deliberate decision to use only free-tier services. The application runs on Next.js 16 deployed to Vercel. The database is MongoDB Atlas (free tier, 512 MB). The AI runs through OpenRouter using four free LLM models in a fallback chain. The entire system costs nothing to host, which matters because most universities that need this kind of tool are the ones that cannot afford commercial alternatives.

The rest of this paper is organized as follows. Section II covers related work on timetabling algorithms and AI-assisted educational tools. Section III describes the system architecture and how the components connect. Section IV explains the CSP-based timetable generation algorithm. Section V covers the RAG-based AI assistant. Section VI presents our experimental setup and results. Section VII discusses what worked, what did not, and what we would do differently. Section VIII concludes.


\section{Related Work}

\subsection{Timetabling algorithms}

The university course timetabling problem (UCTP) has been approached from several directions. Carter and Laporte \cite{carter1996recent} surveyed constraint-based and search-based methods in the mid-1990s, and the basic landscape has not changed as much as you might expect. The problem is still NP-hard, and most practical solutions use some form of heuristic search.

Genetic algorithms have been popular since the early 2000s. Alhadad et al.\ \cite{abdelhalim2016utilizing} used a GA with tournament selection and achieved reasonable schedules for small-to-medium datasets, but convergence was slow (several minutes for 50 subjects). Simulated annealing performs well when the initial solution is already close to feasible \cite{abramson1991constructing}, but it struggles when the constraint space is tight and most random assignments are invalid.

CSP-based methods, which is what we use, treat the problem as variable-domain-constraint triples and apply systematic search with pruning. Qu et al.\ \cite{qu2009survey} reviewed these approaches and noted that backtracking with forward checking works well for medium-sized problems. The MRV (Minimum Remaining Values) heuristic, which schedules the most constrained classes first, consistently reduces the number of backtracks needed \cite{dechter2003constraint}.

More recently, hybrid approaches have combined CSP with metaheuristics. Abdullah et al.\ \cite{abdullah2012hybrid} used a CSP to generate an initial feasible solution and then applied a great deluge algorithm to optimize soft constraints. This two-phase approach is similar to ours, though we use a scoring function during placement rather than a separate optimization pass.

\subsection{AI in educational systems}

Chatbots in education have mostly been used for tutoring and FAQ answering. Wollny et al.\ \cite{wollny2021we} reviewed 53 educational chatbot papers and found that most operate on fixed scripts or simple intent classification. Few connect to live databases, and almost none deal with scheduling data.

RAG (Retrieval-Augmented Generation), introduced by Lewis et al.\ \cite{lewis2020retrieval}, addresses a core weakness of LLMs: they make things up. By retrieving relevant documents before generating a response, RAG grounds the output in actual data. We apply this idea to a structured database rather than a document corpus, which means our retrieval step runs MongoDB queries instead of vector similarity searches.

The use of multiple LLM models with fallback is less studied but practically important. Free-tier API models have rate limits and occasional downtime. Running a single model means the system goes down when the model does. Our four-model fallback chain (StepFun, DeepSeek, GPT-OSS, Qwen) means the chatbot stays available even when individual models are temporarily unreachable.


\section{System Architecture}

The system has four layers: a presentation layer (React components), a routing layer (Next.js API routes), an intelligence layer (CSP generator and AI services), and a data layer (MongoDB). Figure \ref{fig:architecture} shows how these connect.

\begin{figure}[htbp]
\centering
\includegraphics[width=0.48\textwidth]{architecture_diagram.png}
\caption{System architecture. The presentation layer serves five role-specific dashboards. API routes handle authentication and CRUD. The intelligence layer contains the CSP generator and the RAG-based AI assistant. MongoDB stores all persistent data.}
\label{fig:architecture}
\end{figure}

\subsection{Technology stack}

We chose the stack based on two criteria: everything had to be free, and everything had to work together without excessive glue code. Table \ref{tab:stack} lists the components.

\begin{table}[htbp]
\caption{Technology Stack}
\label{tab:stack}
\centering
\begin{tabular}{@{}llp{3.8cm}@{}}
\toprule
\textbf{Layer} & \textbf{Technology} & \textbf{Why we chose it} \\
\midrule
Framework & Next.js 16 & Server and client in one codebase. API routes eliminate the need for a separate backend server. \\
UI & React 19 & Component model fits the multi-dashboard structure. Server components reduce client bundle size. \\
Language & TypeScript 5 & Catches type errors at compile time. The timetable data model is complex enough that untyped code would be painful. \\
Database & MongoDB (Mongoose 9) & Flexible schema. Timetable entries embed naturally as subdocuments. Free Atlas tier gives 512 MB. \\
Auth & NextAuth v5 & Handles sessions, JWT tokens, role-based access. Works with Next.js App Router out of the box. \\
AI & OpenRouter API & Single API endpoint, four free models. No per-token charges. \\
Styling & Tailwind CSS v4 & Utility classes. No custom CSS files to manage. \\
Hosting & Vercel & Zero-config deployment for Next.js. Free tier handles moderate traffic. \\
\bottomrule
\end{tabular}
\end{table}

\subsection{Data models}

The database has nine collections: User, Department, Subject, Room, Batch, TimeSlot, Timetable, Leave, and Notification. The relationships between them are straightforward reference-based links (ObjectId fields), except for timetable entries, which are embedded subdocuments inside the Timetable collection.

A single timetable entry stores: the day (0--5, Monday through Saturday), the time slot index, and references to the subject, faculty member, room, and batch. The entry also records whether it is a theory or practical session. Embedding entries inside the timetable document means we can load an entire schedule in one query, which matters for the AI context-building step.

Users have one of five roles: admin, hod, coordinator, faculty, or student. Each role sees a different dashboard with different capabilities. Admins manage everything. HODs approve timetables for their department. Coordinators create and edit timetables. Faculty view their personal schedules. Students see their batch schedule. The role is stored as a string enum on the User document and checked on every API request through NextAuth session callbacks.

\subsection{Authentication and authorization}

We use NextAuth v5 with a credentials provider. Passwords are hashed with bcryptjs before storage. Sessions use JWT tokens with a 24-hour expiry. The token carries the user's ID, role, and department, so API routes can check permissions without an extra database lookup on every request.

Every API endpoint checks the session first. If the user is not logged in, it returns 401. If they are logged in but their role does not have permission for the requested action, it returns 403. There is no bypass. A student cannot hit the admin API endpoints even if they know the URL.


\section{Timetable Generation Algorithm}

The generator treats scheduling as a CSP. Each class that needs to be scheduled is a variable. The domain of each variable is the set of (day, time-slot, room, faculty) tuples that do not violate any hard constraint. The algorithm places classes one at a time, choosing the most constrained class first (MRV heuristic), and scores each valid placement to prefer better soft-constraint outcomes.

\subsection{Input and preprocessing}

The generator takes as input: a list of subjects (with credit hours, type, and assigned faculty), a list of batches (with student counts), a list of rooms (with capacity and type), and a list of available time slots per day.

The first step is expanding subjects into individual class requirements. A 4-credit theory subject needs 4 lecture slots per week. A 2-credit practical needs 2 lab slots. This expansion produces the full list of classes that need to be placed, typically 150--300 for a single department semester.

\subsection{Constraint model}

\textbf{Hard constraints} are binary: either satisfied or violated. The algorithm never places a class that violates a hard constraint.

\begin{itemize}
\item \textbf{Faculty uniqueness:} A faculty member cannot teach two classes at the same day and slot.
\item \textbf{Room uniqueness:} A room cannot host two classes at the same day and slot.
\item \textbf{Batch uniqueness:} A batch cannot attend two classes at the same day and slot.
\item \textbf{Room type matching:} Practical sessions must go in lab-type rooms. Theory sessions can go in classrooms or lecture halls.
\item \textbf{Room capacity:} The room must seat at least as many students as the batch contains.
\end{itemize}

\textbf{Soft constraints} are scored numerically. The algorithm prefers placements with higher scores but does not reject a placement for violating a soft constraint.

\begin{itemize}
\item \textbf{Workload balance:} Penalizes assigning a faculty member more classes on a day they already have many. Each existing assignment on that day incurs a $-3$ score penalty.
\item \textbf{Consecutive grouping:} Rewards placing lectures of the same subject in adjacent slots ($+15$ bonus).
\item \textbf{Time preference:} Earlier slots get a slight bonus. Morning classes are generally preferred by institutions.
\item \textbf{Faculty preference:} If a faculty member has indicated preferred slots, placements in those slots get $+20$.
\item \textbf{Gap minimization:} Penalizes schedules where a batch has more than 2 free periods between their first and last class of the day.
\item \textbf{Daily limit:} Flags faculty teaching more than 6 hours in a single day.
\end{itemize}

\subsection{Search procedure}

The algorithm uses the following steps:

\begin{enumerate}
\item Sort all class requirements by constraint tightness (MRV). Classes with a specific assigned faculty are scheduled first, then practicals (which need specific room types), then general theory lectures.
\item For each unscheduled class, enumerate all valid (day, slot, room, faculty) combinations. A combination is valid if it violates no hard constraint.
\item Score each valid combination using the soft constraint weights.
\item Select the highest-scoring combination and commit the placement.
\item Update the faculty, room, and batch schedule maps.
\item Repeat until all classes are placed or no valid slot exists for a class.
\end{enumerate}

If a class cannot be placed (step 6), the algorithm records it as unscheduled rather than backtracking, because full backtracking on a greedy placement can be extremely slow and the practical benefit is small. In our tests, fewer than 3\% of classes go unscheduled with this approach, and those are typically cases where the input data itself is infeasible (more classes than available slots).

Figure \ref{fig:flowchart} shows the full generation flow from input to scored output.

\begin{figure}[htbp]
\centering
\includegraphics[width=0.48\textwidth]{algorithm_flowchart.png}
\caption{Flowchart of the CSP-based timetable generation algorithm. Classes are sorted by constraint tightness (MRV), placed in the highest-scoring valid slot, and verified after generation. Classes that cannot be placed are logged for manual assignment.}
\label{fig:flowchart}
\end{figure}

\subsection{Conflict detection and scoring}

After generation, the system runs a verification pass. It checks all hard constraints again (redundantly, as a safety net) and counts soft constraint violations. It also computes an optimization score on a 0--10 scale based on:

\begin{equation}
S = 10 - \frac{H \times 2 + V_{soft}}{N}
\end{equation}

where $H$ is the number of hard constraint violations (should be 0), $V_{soft}$ is the number of soft constraint violations, and $N$ is the total number of placed classes. A score above 7 indicates a good schedule. Below 5 suggests the input data is too constrained and the coordinator should add more rooms or adjust faculty loads.

The conflict detector also runs independently when coordinators manually edit a timetable. If they swap a room or move a class to a different slot, the system immediately checks whether the change introduces any conflicts and warns them before saving.


\section{RAG-Based AI Assistant}

The AI assistant answers natural language questions about schedules using data from the live database. It is not a general-purpose chatbot. It does not answer questions about the weather or write poetry. If you ask it something unrelated to the institution's timetable data, it says so.

\subsection{Query processing pipeline}

When a user sends a message, the system processes it in four steps:

\textbf{Step 1: Keyword extraction.} The system parses the query using regex patterns to identify: day names (Monday, Tuesday, etc.), room identifiers (alphanumeric patterns like "R101" or "Lab 3"), time expressions ("10 AM", "morning"), and subject codes. It also strips stopwords and extracts general keywords for broader matching.

\textbf{Step 2: Intent classification.} Based on the extracted keywords, the system classifies the query into one of seven intents: schedule lookup, faculty query, room query, subject query, conflict query, availability query, or statistics query. The classification is rule-based (keyword matching), not ML-based. This is intentional. Intent classification for a narrow domain with 7 categories does not need a neural network; pattern matching is faster and more predictable.

\textbf{Step 3: Context retrieval.} Based on the classified intent, the system runs targeted MongoDB queries. A schedule query fetches timetable entries matching the identified day, batch, or faculty. A room query fetches room details and their current bookings. A statistics query aggregates counts across collections. Each retrieved record gets a relevance score. The top results (up to a configurable limit) are formatted into structured text blocks.

\textbf{Step 4: Prompt construction and LLM call.} The retrieved context, along with system-level statistics (total counts of departments, subjects, rooms, etc.) and the conversation history (last 5 messages), is assembled into a prompt. The prompt includes a system message that instructs the model to answer based only on the provided data and to say "I don't have that information" rather than guessing.

\subsection{Multi-model fallback}

The system connects to OpenRouter, which provides a unified API for multiple LLM providers. We configured four free models:

\begin{enumerate}
\item StepFun 3.5 Flash (primary)
\item DeepSeek R1
\item GPT-OSS 120B
\item Qwen3 Coder
\end{enumerate}

Each model has its own API key (OpenRouter allows multiple free-tier keys). When a query comes in, the system tries Model 1 first. If it fails (timeout, rate limit, server error), it tries Model 2, then Model 3, then Model 4. If all four fail, it returns a static error message. In practice, we have never seen all four fail simultaneously. Usually one or two are down at any given time, and the fallback handles it transparently.

The response includes which model answered, so users (and developers during debugging) know where the answer came from. Response times vary by model: StepFun is typically fastest (1--2 seconds), DeepSeek is slowest but most detailed (3--5 seconds).

\subsection{Specialized AI services}

Beyond general Q\&A, the system offers two specialized AI modes:

\textbf{Conflict resolution.} When the timetable has conflicts (detected by the verification pass or introduced by manual edits), the AI can analyze them. It receives the specific conflict details (which room/faculty/batch, which time slots) and generates an analysis: root cause, 3--5 suggested fixes, and prevention advice. The prompt is tailored to produce actionable output rather than generic suggestions.

\textbf{Timetable optimization.} The AI can review a complete timetable and suggest improvements. It receives the full schedule, the optimization score, the constraint violation counts, and the room/faculty utilization rates. It then produces a prioritized list of changes that could improve the score.

Both modes use RAG context to ground their suggestions in actual data. If the AI suggests moving a class to Room 301, it has already checked (via the retrieval step) that Room 301 exists and has the right capacity.

\subsection{Quick FAQ fallback}

For common questions ("How do I generate a timetable?", "What are the user roles?", "How do I add a room?"), the system returns hardcoded answers without calling any LLM. This reduces latency for frequent queries from 2--5 seconds to under 100 milliseconds and avoids wasting API calls on questions that always have the same answer.


\section{Experimental Evaluation}

\subsection{Test environment}

We tested on a laptop (Intel i5-12400, 16 GB RAM, Windows 11) running the application locally in development mode (next dev). The database was hosted on MongoDB Atlas free tier (M0 cluster, shared instance, US-East-1). AI queries went through the public OpenRouter API.

This is not a high-performance test environment, and that is the point. If the system works on a mid-range laptop with free-tier cloud services, it will work in the environments where it is likely to be deployed.

\subsection{Dataset}

We created synthetic test data representing a mid-sized university department:

\begin{itemize}
\item 6 departments (CS, IT, ECE, EEE, ME, CE)
\item 40 faculty members (distributed across departments)
\item 120 subjects (theory and practical, 2--4 credits each)
\item 30 rooms (20 classrooms, 10 labs, capacity 30--120)
\item 20 batches (4 years $\times$ 5 departments, ~60 students each)
\item 6 time slots per day, 6 days per week (Monday--Saturday)
\end{itemize}

\subsection{Timetable generation results}

We ran the generator 20 times for each department (120 runs total) and recorded the results. Table \ref{tab:generation} summarizes the performance.

\begin{table}[htbp]
\caption{Timetable Generation Performance Across 120 Runs}
\label{tab:generation}
\centering
\begin{tabular}{@{}lcccc@{}}
\toprule
\textbf{Metric} & \textbf{Min} & \textbf{Max} & \textbf{Mean} & \textbf{Std Dev} \\
\midrule
Generation time (s) & 4.2 & 28.7 & 12.3 & 5.1 \\
Classes placed (\%) & 94.1 & 100.0 & 97.8 & 1.6 \\
Hard violations & 0 & 0 & 0 & 0 \\
Soft violations & 2 & 18 & 7.4 & 3.2 \\
Optimization score & 6.1 & 9.4 & 7.8 & 0.8 \\
\bottomrule
\end{tabular}
\end{table}

Every run produced zero hard constraint violations. The generator never double-booked a room, faculty member, or batch. Soft violations averaged 7.4 per schedule, mostly gaps in student schedules (batches with free periods between classes). The optimization score averaged 7.8 out of 10, which indicates reasonably good schedules.

The 2.2\% of unscheduled classes (across all runs) were almost entirely practicals for large batches that needed labs with 100+ capacity. We only had 2 labs of that size in our test data. Adding one more large lab to the dataset reduced unscheduled classes to under 0.5\%.

Generation time depended mostly on the number of subjects and the constraint tightness. Departments with many lab courses took longer because the room-type constraint narrows the domain significantly. The 28.7-second worst case was a deliberately adversarial configuration with 25 subjects competing for 3 lab rooms.

\subsection{AI assistant evaluation}

We tested the chatbot with 50 queries across the seven intent categories. We manually judged each response for:

\begin{itemize}
\item \textbf{Accuracy:} Does the response contain correct information?
\item \textbf{Groundedness:} Is the information actually from the database, not hallucinated?
\item \textbf{Relevance:} Does it answer what was asked?
\item \textbf{Latency:} How long did it take?
\end{itemize}

Accuracy was 88\% (44/50). The 6 incorrect responses were edge cases: queries mentioning subjects by informal names not in the database ("DBMS" instead of "Database Management Systems"), and one case where the LLM ignored the context and answered from its training data instead.

Groundedness was 92\% (46/50). Four responses included statements not backed by the retrieved context, though they happened to be correct (general knowledge about university scheduling). This is a known RAG weakness: models sometimes blend retrieved context with parametric knowledge.

Relevance was 96\% (48/50). Two responses went off-topic, both when the query was ambiguous ("tell me about Monday" without specifying which batch or faculty).

Average latency was 2.8 seconds. The breakdown: keyword extraction and intent classification took under 50 ms, MongoDB queries took 200--500 ms, and the LLM call took 1.5--4 seconds depending on which model responded. Quick FAQ responses averaged 80 ms.

The fallback mechanism triggered 7 times out of 50 queries (14\%). StepFun was unavailable for about 20 minutes during our testing window, and the system silently switched to DeepSeek. Users would not have noticed the switch except that responses took slightly longer.

\subsection{Comparison with manual scheduling}

We asked two faculty coordinators from our department to manually create a timetable for the same input data. They used Google Sheets (their usual method). Coordinator A took 3 hours and 40 minutes and produced a schedule with 4 hard conflicts (a room double-booking and 3 faculty overlaps) that took another 45 minutes to fix. Coordinator B took 2 hours and 15 minutes with 1 hard conflict, resolved in 10 minutes.

Our system produced a conflict-free schedule in 12 seconds. This is not entirely a fair comparison---the coordinators were also dealing with preferences and politics that our system does not model---but the time difference is real, and the conflict-free guarantee matters.


\section{Discussion}

\subsection{What works well}

The CSP generator is reliable. In 120 test runs, it never produced a hard constraint violation. The MRV heuristic does most of the heavy lifting---by scheduling the most constrained classes first, the algorithm avoids painting itself into corners. This is a known result in CSP literature \cite{dechter2003constraint}, but it is satisfying to see it work in practice.

The multi-model fallback makes the AI assistant much more robust than depending on a single model. Free-tier LLMs have unpredictable downtime. Having four options means the chatbot stayed available through every test session we ran.

The zero-cost deployment is probably the most practically important aspect. We have seen university departments try to use scheduling software and abandon it because of licensing fees. A system that runs entirely on free tiers removes that barrier.

\subsection{What could be better}

The keyword-based RAG retrieval is the weakest link. It works for straightforward queries ("Show me Room 301 schedule on Monday") but fails when users phrase things in unexpected ways. A vector-embedding approach would handle semantic similarity better, but it would add cost (embedding API calls or a local embedding model) and complexity.

The greedy CSP solver without full backtracking means the algorithm occasionally leaves classes unscheduled that could be placed with a more exhaustive search. We chose speed over completeness. For the rare cases where a class cannot be placed, the coordinator can manually assign it or adjust constraints and regenerate.

The optimization scoring is basic. The linear formula in Equation (1) treats all soft violations equally, but in practice, a faculty member teaching 7 hours in one day is worse than a batch having 3 gaps. A weighted scoring function would produce better rankings but would need per-institution tuning.

\subsection{Limitations}

The system does not handle multi-section courses (where the same subject has multiple sections with different faculty). It treats each subject-batch pair as a single stream. Supporting sections would require extending the data model and the generator.

The AI assistant cannot modify the timetable. It can analyze and suggest, but acting on suggestions requires a human going to the relevant page and making the change. An agentic AI that can directly execute scheduling modifications is a possible future direction, but it introduces safety concerns (you probably want a human in the loop for actual schedule changes).

We tested with synthetic data, not a real institution's data. Real data has messier constraints: faculty who teach across departments, shared labs, time slots blocked for meetings, exam periods. The system supports these through the constraint model, but we have not validated it with real-world edge cases.


\section{Conclusion}

We built a timetable scheduling system that generates conflict-free schedules using a CSP solver and lets users query those schedules through an AI chatbot. The CSP engine uses MRV ordering and greedy placement with soft-constraint scoring. The chatbot uses keyword-based RAG to ground responses in actual database records and falls back across four free LLM models to stay available.

The system works. It generates reasonable schedules fast, answers most questions correctly, and costs nothing to deploy. That last point is not a footnote---it is the whole reason we built it this way. Commercial scheduling software exists. Commercial AI APIs exist. But a department that cannot afford them is stuck with spreadsheets and frustrated coordinators. This system gives them an alternative.

The code is open-source and designed to run on free-tier infrastructure (Vercel, MongoDB Atlas, OpenRouter). It handles five user roles, CRUD operations for all entities, and a conversational AI that does not hallucinate room numbers. There is room for improvement---vector embeddings for RAG, weighted optimization scoring, multi-section support---but the current version solves the core problem well enough to be useful.

If you are a university coordinator reading this and you are tired of the spreadsheet, the code is on GitHub. The README has setup instructions. It takes about 15 minutes.

\section*{Acknowledgments}

We thank our project guide for their feedback throughout the development process. We also acknowledge the open-source communities behind Next.js, MongoDB, and the free-tier LLM providers (StepFun, DeepSeek, Alibaba Qwen, OpenAI GPT-OSS) who make projects like this possible without funding.

\begin{thebibliography}{00}

\bibitem{even1975complexity} S.\ Even, A.\ Itai, and A.\ Shamir, "On the complexity of timetable and multicommodity flow problems," \textit{SIAM Journal on Computing}, vol.\ 5, no.\ 4, pp.\ 691--703, 1976.

\bibitem{cooper1996complexity} T.\ B.\ Cooper and J.\ H.\ Kingston, "The complexity of timetable construction problems," in \textit{Practice and Theory of Automated Timetabling}, Lecture Notes in Computer Science, vol.\ 1153, Springer, 1996, pp.\ 281--295.

\bibitem{carter1996recent} M.\ W.\ Carter and G.\ Laporte, "Recent developments in practical course timetabling," in \textit{Practice and Theory of Automated Timetabling}, Lecture Notes in Computer Science, vol.\ 1153, Springer, 1996, pp.\ 3--19.

\bibitem{abdelhalim2016utilizing} E.\ A.\ Abdelhalim and G.\ A.\ El Khayat, "Utilizing genetic algorithms for solving university course timetabling problem," \textit{Journal of Intelligent Learning Systems and Applications}, vol.\ 8, no.\ 2, pp.\ 29--41, 2016.

\bibitem{abramson1991constructing} D.\ Abramson, "Constructing school timetables using simulated annealing: sequential and parallel algorithms," \textit{Management Science}, vol.\ 37, no.\ 1, pp.\ 98--113, 1991.

\bibitem{qu2009survey} R.\ Qu, E.\ K.\ Burke, B.\ McCollum, L.\ T.\ G.\ Merlot, and S.\ Y.\ Lee, "A survey of search methodologies and automated system development for examination timetabling," \textit{Journal of Scheduling}, vol.\ 12, no.\ 1, pp.\ 55--89, 2009.

\bibitem{dechter2003constraint} R.\ Dechter, \textit{Constraint Processing}. Morgan Kaufmann, 2003.

\bibitem{abdullah2012hybrid} S.\ Abdullah, H.\ Turabieh, B.\ McCollum, and P.\ McMullan, "A hybrid metaheuristic approach to the university course timetabling problem," \textit{Journal of Heuristics}, vol.\ 18, no.\ 1, pp.\ 1--23, 2012.

\bibitem{wollny2021we} S.\ Wollny, J.\ Oliveira, W.\ Streuber, R.\ Bloch, H.\ Lujic, and J.\ Grover, "Are we there yet? A systematic literature review on chatbots in education," \textit{Frontiers in Artificial Intelligence}, vol.\ 4, p.\ 654924, 2021.

\bibitem{lewis2020retrieval} P.\ Lewis, E.\ Perez, A.\ Piktus, F.\ Petroni, V.\ Karpukhin, N.\ Goyal, H.\ Kuttler, M.\ Lewis, W.\ Yih, T.\ Rocktaschel, S.\ Riedel, and D.\ Kiela, "Retrieval-augmented generation for knowledge-intensive NLP tasks," in \textit{Advances in Neural Information Processing Systems}, vol.\ 33, 2020, pp.\ 9459--9474.

\bibitem{nextjs2024} Vercel, "Next.js Documentation," 2024. [Online]. Available: https://nextjs.org/docs

\bibitem{mongoose2024} Mongoose, "Mongoose ODM Documentation," 2024. [Online]. Available: https://mongoosejs.com/docs/

\end{thebibliography}

\end{document}
